%
% File main.tex
%
% NLP course project-proposal template, Warsaw University of Technology.
% Prepared by Anna Wróblewska, adapted from the ACL 2015 conference
% formatting template (acl2015.sty), originally by the ACL 2015
% organizers and earlier ACL/EACL format authors (see acl2015.pdf
% for the full credits chain).

\documentclass[11pt]{article}
\usepackage{acl2015}
\usepackage{times}
\usepackage{url}
\usepackage{latexsym}
\usepackage{xcolor}

%\setlength\titlebox{5cm}

% You can expand the title box if you need extra space
% to show all the authors. Please do not make the title box
% smaller than 5cm (the original size); we will check this
% in the camera-ready version and ask you to change it back.


\title{21. When Does a RAG Framework Actually Matter? A Systematic Review
and Controlled Experimental Comparison
\\Project Proposal for NLP Course, Winter 2026/27}
% https://www.overleaf.com/project/63440c16442dba647b52a1ae
\author{Jan Czemko (submitted by) \\
  Album number \\
  {\tt email@pw.edu.pl} \\\And
  Karolina Móżdżyńska \\
  Album number \\
  {\tt email@pw.edu.pl} \\ \And
  Anna Piórkowska \\
  354873 \\
  {\tt anna.piorkowska3.stud@pw.edu.pl} \\ \AND
  supervisor: Anna Wróblewska\\
  Warsaw University of Technology \\
  {\tt anna.wroblewska1@pw.edu.pl}}
\date{}

\begin{document}
\maketitle

{\small\em Template: Anna Wróblewska (Warsaw University of Technology), adapted from the ACL 2015 conference formatting template.}
\vspace{4pt}

\begin{abstract}
  Write about the proposed work touching the research gap you are addressing, novelty, and on the proposed deliverable aimed from this project.
  (Short description up to 250 words.)
\end{abstract}
% w przegladzie literatury uwzglednilismy flashraga, longraga, miniraga i na podstawie dokumentacji langchain i llamaindex
\section{Introduction}
Retrieval-augmented generation (RAG) is a technique that helps large language models to retrieve and incorporate new information from external data sources, that supplement knowledge from the LLM's pre-existing training data. 
% albo inne nazwy frameworków które będziemy badać
Many frameworks such as LangChain, LlamaIndex, Haystack, offer seemingly similar capabilities. They have different level of abstraction, default pipelines, component choices. Benchmarks comparing some frameworks usually changes it together with embedding model, retrieved context, prompts and some other parameters. Then it is unclear what caused the performance to differ, the framework itself, or its default settings. The project tests how different retrieval and answer quality of popular RAG frameworks is, when everything except their models, context, embeddings and prompts held constant. \\
\section{Data}
Question-answering systems make it easier to find answers when someone is searching for information. However, developing such tools requires appropriate test datasets that accurately reflect the complexity of the problems involved. Many existing datasets for question-based information retrieval are limited to simple, factual queries. Quasper is a dataset designed to bridge this gap, it contains 5,049 questions based on 1,585 scientific papers in the field of natural language processing. The questions and answers were prepared by individuals with expertise in NLP. First, they read the title and abstract of a paper and formulated a question concerning the work as a whole. Subsequently, a different group provided answers, complete with justifications and references to specific excerpts from the text. The results indicate that models performing well on other datasets achieve significantly lower scores on this one. Consequently, this dataset serves as an effective tool for evaluating and comparing various RAG (Retrieval-Augmented Generation) frameworks. The data includes the `id ` and the following labels: `title`, and `abstract` (of type string) and `full\_text`, `qas` (of type sequence). Each document contains an average of 3.2 questions. The `qasper-dev-v0.1.json` file, which contains 281 rows, will be used as a benchmark set. It is licensed under CC BY 4.0 and was obtained from Hugging Face.\\

HotpotQA is a dataset containing 113,000 question-answer pairs that require multi-hop reasoning, based on Wikipedia. It was created by an NLP research team. The dataset includes the following labels: `answer`, `context`, `id`, `level`, `question`, `supporting\_facts`, and `type`. From the "distractors" subset, we select 150 questions for the fixed test set. The HotpotQA dataset is distributed under the CC BY-SA 4.0 license, and the data was obtained from the Hugging Face platform.
\\Detailed description (\textbf{max. 4 meaningful pages}) in English, including the following information:
\begin{enumerate}
\item scientific goal of the project (description of the problem to be solved, research questions and hypotheses);
\item significance of the project (state of the art, justification for tackling a specific scientific problem, justification for the pioneering nature of the project, the impact of the project results on the development of the research field and scientific discipline);
\item concept and work plan (general work plan - with a timeline, specific research goals, results of preliminary research, risk analysis);
\item approach \& research methodology (underlying scientific methodology, methods, techniques and research tools, methods of results analysis, equipment and devices to be used in research);
\item well-justified project literature (a reference list for publications included in the project description, with full bibliographic data);
\item \textbf{data statement} (new for 2026/27) -- source of the data, licence, whether it contains personal data, and what you do about it. Six lines are enough;
\item \textbf{contribution table} -- each team member, what they did, workload in hours. \textbf{Required in this and every later phase report; a report submitted without it is not accepted for grading, not just marked down};
\item \textbf{declaration of LLM use} -- which tools, for what. \textbf{A report without this is not accepted for grading} (new for 2026/27).
\end{enumerate}

\textbf{Reminder: place this and every later submission in a folder named \texttt{surname\_surname\_surname\_surname} (all team members). The order set here must not change for the rest of the semester -- a mismatch costs 5 points.}


\section{References}


% include your bib file like this:
%\bibliographystyle{acl}
%\bibliography{acl2015}

\end{document}